\chapter{Formal Proof with SPARK}
\label{ch:spark}\index{SPARK}\index{formal proof}\index{proof!SPARK}

Every layer in Chapter~\ref{ch:testing} answers the same question in a different
way: \emph{does the code do the right thing on the inputs we tried?} A test runs one
input and checks one answer; a thousand tests run a thousand. Formal proof asks a
different question, and answers it once for all inputs: \emph{can this code go wrong
on \textbf{any} input at all?} Where a test demonstrates the presence of correct
behaviour on an example, a proof demonstrates the \emph{absence} of a whole class of
failure over the entire input space. That is the layer this chapter adds, and for a
bare-metal system that parses bytes it never chose --- a certificate from a stranger,
a GPS sentence from a noisy sky, a DNS reply from the network --- it is the layer
that matters most.

SPARK is a subset of Ada with a proof tool, \texttt{gnatprove}, that ships in the
same Alire toolchain as the compiler. Mark a package \lstinline{with SPARK_Mode => On}
and the tool reads its contracts and its code and discharges, or fails to discharge,
a \emph{verification condition} for every operation that could fault: every array
index, every arithmetic operation that could overflow, every division, every
dereference, and the termination of every loop. When it succeeds it has proved, with
an SMT solver rather than a test, that the code is free of run-time errors --- the
\emph{silver} level of SPARK adoption, absence of run-time errors (\index{AoRTE}AoRTE).
No \lstinline{Constraint_Error}, no \lstinline{Storage_Error} from an unbounded loop,
no buffer overrun, on any input the types allow.

\section{What is proved}\index{proof!surface}

The proof surface is the pure, bounded logic of the HAL: the parsers, serializers,
checksums, and conversions that compute over bytes rather than poke registers. Around
thirty units carry \lstinline{SPARK_Mode => On} and discharge some seven hundred and
eighty run-time checks with zero unproved. The set falls into natural groups: the
ext4 on-disk serializers and their byte helpers (Chapter~\ref{ch:storage}); the
network and protocol parsers; the software half of the cryptography; and a scattering
of pure arithmetic --- civil-date math, longest-prefix routing, sensor conversions,
and the timing dividers extracted from the PWM drivers.

The headline is the group exposed to hostile input. The X.509 certificate reader, the
NMEA GPS-sentence parser, the DNS response parser, the certificate-chain walk, and the
Modbus slave and master all parse bytes chosen by someone else, and all are now proved
to have \textbf{no buffer overrun, no overflow, and no infinite loop on any malformed
or malicious message}. This is not a hardened version that passed a fuzzer; it is a
proof that the failure cannot occur. The DNS parser makes the point sharply: a DNS
name can contain a \emph{compression pointer} back into the message, and a malicious
pointer that points at itself would loop a naive parser forever. The proof required a
\lstinline{Loop_Variant} showing the cursor strictly advances, so termination on a
self-referential pointer is not tested but guaranteed.

\section{Drawing the boundary}\index{SPARK!boundary}

A driver is not pure. It writes volatile registers, follows access-type pointers, and
raises exceptions on bad hardware --- none of which is in the SPARK subset. The
technique is to draw a line \emph{inside} the unit. The pure computation is marked
\lstinline{On} and proved; the parts that touch the world are marked
\lstinline{SPARK_Mode => Off} and left to the other verification layers. For the ext4
serializers this meant \emph{factoring}: the buffer-to-record step was lifted out of
each \lstinline{Read}/\lstinline{Write} into a pure \lstinline{Decode}/\lstinline{Encode}
function that a proved caller invokes, while the block-device I/O and the checksum
stayed behind the boundary. The move is behaviour-neutral --- the same expressions,
relocated --- and the host filesystem suite (Chapter~\ref{ch:testing}, Layer~2) reran
\texttt{e2fsck}-clean afterward to confirm it.

One recurring obstacle is worth recording because it shapes every buffer proof. An
\lstinline{array (Natural range <>)} of bytes can, \emph{in the type}, have a length
of $2^{31}$, which overflows \lstinline{Integer}; so the solver cannot even prove that
\lstinline{Off + N <= Buf'Length} does not overflow, and the whole parser stalls. The
fix is to constrain the index --- \lstinline{subtype Buffer_Index is Natural range 0 .. 2**24 - 1}
--- so the length provably fits an integer. A certificate or a filesystem block is far
below sixteen megabytes, so the constraint costs nothing and unlocks everything; it is
the single change that made the ext4 and X.509 buffers provable.

\section{What the proofs found}\index{bug!found by proof}

A proof that goes through is quiet. A proof that will not go through is often pointing
at a real defect, and this effort surfaced six. Five were on a malformed-input path
that testing had not reached; the sixth was a different animal, and worth reading
last:

\begin{itemize}
  \item \textbf{Modbus \lstinline{Put_MBAP}} computed \lstinline{PDU_Len + 1} on an
    unbounded length --- an overflow on a crafted request. The proof drove the
    precondition \lstinline{PDU_Len <= Max_PDU}, the real 253-byte protocol cap.
  \item \textbf{Superblock \lstinline{Encode}} checksummed an \emph{absolute}
    \lstinline{Buf (Base .. Base + N)} slice where every writer used
    \lstinline{Buf'First + offset} --- wrong for a non-zero-based buffer; now
    \lstinline{Buf'First}-relative.
  \item \textbf{X.509 \lstinline{Host_Matches}} could index out of range on a hostile
    certificate whose subject-alt-name slices did not match its buffer; now guarded.
  \item \textbf{Chain validation} dereferenced a certificate's \lstinline{Data} pointer
    with no null or empty check --- a caller-supplied chain could crash it; now rejected
    as malformed.
  \item \textbf{The SHT41 CRC} loop bound \lstinline{while I + 2 <= Data'Last} overflowed
    in its own guard and underflowed on an empty buffer; rewritten as a group count.
  \item \textbf{\lstinline{ADC.Read} returned a stale conversion.} SPARK refused it as
    a ``function with output global'': it started a conversion --- writing the SAR
    registers --- and recorded the result's validity in a package-level flag. The
    conversion-done poll has a deadline, and on expiry the data register still holds
    whatever it last latched, so the function returned that as though it were fresh.
    The only evidence was a global that a second task could overwrite in between. It is
    now a procedure reporting validity per call.
\end{itemize}

The first five were not reachable by a well-formed message, which is exactly why the
test suites had not caught them and why proving over the whole input space did.

The sixth is worth separating, because it was not an arithmetic slip on a hostile input
and no amount of fuzzing would have found it: the \emph{signature was wrong}. A
function that starts a conversion is not a function, and SPARK objects to the shape
before it ever considers the values. That objection recurred four times in the
application this driver was written for --- a current reading, a button poll, a fault
flag, a console token cursor --- each a question that changed the world on the way to
answering. Nothing about them looked alike in the source. The tool named them with the
same sentence every time, which is a fair description of what a prover is for: not
cleverness, but never tiring of the same question.

\section{Proving code that touches registers}\index{SPARK!registers}\label{spark:registers}

A memory-mapped register is not an ordinary variable, and until the generated register
layer said so, no unit that touched one could be analysed. \texttt{svd2ada} marks each
peripheral record \lstinline{Volatile}, which is enough for the compiler and not enough
for a prover: with the four volatile aspects unstated they all default to
\lstinline{True}, the most conservative reading, in which \emph{a read is itself a state
change}. Under that default SPARK rejects any \lstinline{Volatile_Function} over a bank
outright --- and a volatile function is how a hardware reading gets into a contract.

The fix is to say which kind of external interaction each bank is:

\begin{lstlisting}
GPIO_Periph : aliased GPIO_Peripheral
with Import, Address => GPIO_Base, Volatile,
     Async_Readers    => True,
     Async_Writers    => True,
     Effective_Reads  => False,   --  reading a pad does not consume it
     Effective_Writes => True;
\end{lstlisting}

The aspects apply to the \emph{object}, and a peripheral is one object, so a bank gets
one characterisation covering all its registers. That forces a choice on
\lstinline{Effective_Reads}, and the two directions are not symmetric.
\lstinline{False} on a bank whose reads consume is \emph{unsound} --- the prover would
believe that reading a FIFO twice yields the same byte. \lstinline{True} on a bank whose
reads do not consume is merely pessimistic. So the banks with a consuming data port keep
\lstinline{True}: \texttt{I2C0/1}, \texttt{RMT}, \texttt{SDHOST}, \texttt{UART0/1/2} and
\texttt{USB\_DEVICE}, chosen by reading the generated registers rather than by name.
\texttt{DMA} and \texttt{TWAI0} look like FIFOs and are not: their counters and buffer
windows do not advance on a read.

The aspects are inert for code generation --- an application binary was byte-identical
before and after --- and they are applied by a post-process in
\texttt{regenerate.sh} rather than by hand, since line~1 of that script is
\lstinline{rm -rf $HERE/svd}.

\textbf{What it buys, and what it costs.} \lstinline{ESP32S3.GPIO}'s lock-free pad
operations --- \lstinline{Set}, \lstinline{Clear}, \lstinline{Write}, \lstinline{Read}
--- are now proved. The substance is not the count but one property: pads
\texttt{0\,..\,31} and \texttt{32\,..\,48} live in different register banks, and no path
can reach one bank holding a pin belonging to the other. \lstinline{Configure} and
\lstinline{Toggle} stay outside, because both serialise a read-modify-write through a
protected object.

The cost is a real constraint on how such code is written. A volatile read may not sit
inside a larger expression, so

\begin{lstlisting}
return (Reg.GPIO_Periph.IN_k and Lo_Bit (Pin)) /= 0;   --  rejected
\end{lstlisting}

\noindent becomes a read into a local, then the mask. That is not pedantry: the pad is
driven from outside the program, and the value must be fixed before anything else looks
at it.

Two further rules bite in practice. Contracts must be \emph{per-subprogram} rather than
package-wide, because a package-wide \lstinline{SPARK_Mode} puts the driver's protected
object inside a SPARK region and fails with ``tasking in SPARK requires sequential
elaboration'' --- which would impose that policy on every client. And the banks must be
named \emph{directly} in \lstinline{Global}: an abstract state cannot work, because a
\lstinline{Refined_State} may only name state declared in the same package, and these
banks belong to the generated layer and are shared with every other driver.

\section{The edge of the subset}\index{SPARK!limits}

Proof has an honest boundary, and naming it is part of the discipline. Three things sit
outside it. First, code that is \emph{out of the SPARK subset by construction}: the
hardware crypto accelerators (SHA, AES, RSA), the controlled-type bus sessions with
their \lstinline{Finalize}, the access-to-subprogram interrupt callbacks. These are
verified by the other layers --- loopback self-tests and known-answer tests --- not by
proof, and that is the right tool for them.

This list used to begin with ``the volatile memory-mapped register drivers'', and that
was wrong. It is left recorded rather than quietly removed, because the mistake is
instructive: a tool's refusal had been read as a statement about what \emph{can} be
proved, when it was a statement about what the code had \emph{told} it. Section
\ref{spark:registers} is the correction. Second, pure arithmetic that is \emph{inlined} in a register procedure, like a
baud prescaler or a dead-time divider; that logic \emph{can} be proved, but only after
extracting it into a pure sibling package, which is exactly what was done for the
TWAI, LEDC, RMT, and MCPWM timing math. Third, code whose \emph{shape} is wrong rather than its content: a SPARK function may not
have out parameters, which is why \lstinline{Public_Key}, \lstinline{ECDH} and
\lstinline{Sign} sit outside --- and in \lstinline{Sign}'s case not the only why, since its
body also calls SPARKNaCl's HMAC-SHA-256, which puts it in the first category as well.
Knowing which of the three a unit falls into is what keeps the surface honest.

That third category used to name something else, and the replacement is the more
instructive story. It read, in full: ``the P-256 signature check is legal SPARK and the
tool runs, but proving it pulls in the entire vendored SPARKNaCl elliptic-curve contract
closure, which did not converge under a long budget --- a dedicated, lemma-level effort
rather than a harvest.'' Two claims, and both were wrong. SPARKNaCl was not being pulled
in: it is named only by the RFC-6979 signing glue, which is outside the subset, and the
proof never reaches it. And what defeated the analysis was not elliptic-curve mathematics
but \emph{inlining}: GNATprove inlines a call to a local subprogram that carries no
contract, and \lstinline{Verify} calls \lstinline{Scalar_Mul} twice, which loops 256 times
over \lstinline{Dbl} and \lstinline{Add}, which call \lstinline{Mont_Mul}, whose CIOS body
is itself two nested loops. Unfolded into one verification condition that is hopeless, and
no amount of lemma work on the modular arithmetic would have helped, because the
arithmetic was never the obstacle. A wrong diagnosis is worth quoting exactly, because the
cost of this one was not the effort it predicted --- it was the years it would have
deferred.

The fix is to stop the unfolding, and there are two levers for that. The blunt one is
\texttt{gnatprove --no-inlining}, which switches contextual analysis off for the whole run:
on the untouched source, with nothing but the two \lstinline{SPARK_Mode => Off} aspects
removed, that alone proves the unit --- 95 obligations, none unproved, in 2m34s. It is
worth knowing that a single switch was ever all it took. The lever this SDK actually
ships is the other one: any contract at all makes a call opaque, so each primitive got
\lstinline{Global => null} --- a statement that is simply true, since they are pure
functions of their arguments, and one SPARK checks rather than takes on trust. That costs
more obligations than the switch (144 against 95) and buys durability: a contract travels
with the source, so the unit proves under any project and any invocation, including the
cross \lstinline{tls.gpr}, while a switch has to be remembered by whoever runs the tool
next.

Either way \lstinline{Verify} and \lstinline{On_Curve} are now proved, and the unit proves
in \emph{less} time than it did with them excluded altogether: 144 obligations in 2m33s,
against 210 in 3m03s before, because inlining had been re-proving the same callee body at
every call site. Two cautions on reading that number. It is \emph{absence of run-time
errors} and nothing more --- the proof says \lstinline{Verify} cannot overflow or run off
an array on any input an attacker chooses, not that it accepts exactly the valid
signatures, which is what the known-answer test is for. And of the contracts added, only
the \lstinline{Global} aspects do the proof any work: the pre- and postconditions on
\lstinline{Inv32}, \lstinline{Inv_Mod}, \lstinline{Is_Zero}, \lstinline{"="} and
\lstinline{Geq} discharge nothing in \lstinline{Verify} and were kept as checked
documentation, not as load-bearing structure. The lesson generalises past this unit ---
when a proof will not converge, ask what the prover was handed before assuming the
difficulty is real.

\section{Running the proofs}\index{book!prove}

The proofs run the way the tests do, from one script. \texttt{book/prove/prove.sh}
invokes \texttt{gnatprove} at the silver level over each native prove project and exits
non-zero if any run-time check is unproved, so a change that breaks a proof breaks the
build just as a failing test would. The proofs run host-native because run-time-error
freedom of pure logic is target-independent --- the same argument as the host tests,
one layer up. The companion \texttt{book/prove/README.md} lists the proved units, the
factoring pattern, and the toolchain gotchas learned along the way. Formal proof does
not replace the layers below it; it caps them. The compiler rejects the ill-typed, the
host tests check the answers, the self-tests check the wiring, the sweep checks the
foundation --- and above them all, proof rules out an entire category of failure on
inputs no test will ever enumerate.
